% Independent single-column Figure 2; native width 3.34 inches.
% Static estimates only. Reproduce via experiments/data/reuse-aiconfigurator-multifamily-a100/.
\begin{figure}[t]
  \centering
  \includegraphics[width=\linewidth]{figures/motivation-row/motivation-recompute.pdf}
  \caption{Estimated recovery latency by history placement.
    Color denotes the recovery path; line style and marker denote
    Qwen3-4B, Llama-3.1-8B, and Qwen3-14B on shared axes.
    A100-SXM, BF16, batch one, TP one; K~=~1,024. All paths include
    the same known 64-token tail. HBM holds the entire history;
    DRAM adds serialized H2D at an assumed 32\,GB/s, without overlap.
    All curves use AIConfigurator HYBRID and a transfer model, not
    physical recovery measurements. Loading, detection, queueing,
    pre-failure replication, and software takeover overheads are excluded.
    Beyond-config context lengths are performance shapes, not accuracy
    validation.}
  \label{fig:motivation-recompute}
  \Description{One plot compares nine estimated recovery curves on shared
    logarithmic context-length and latency axes. Blue, orange, and red
    encode HBM reuse, DRAM reuse, and empty-cache reconstruction.
    Circles with solid lines, squares with dashed lines, and triangles
    with dash-dot lines encode Qwen3-4B, Llama-3.1-8B, and Qwen3-14B,
    respectively.
    All paths rebuild the same known 64-token tail.
    At 128K history tokens, the three models span 35 to 59 milliseconds
    for HBM, 572 to 730 milliseconds for DRAM, and 31.4 to 49.6 seconds
    for Recompute. DRAM transfer bandwidth is assumed, not measured.}
\end{figure}
